%% -*- coding: utf-8 -*-
%% EXAMPLE CONTENT — ships with a new SciTeX workspace. Replace it with yours.
\section{Methods}
\label{sec:methods}
\pdfbookmark[1]{Methods}{methods}

\subsection{Data}

The images come from the UCI \emph{Optical Recognition of Handwritten Digits}
dataset of Alpaydin and Kaynak, in which 43 contributors wrote digits that were
captured as $32\times32$ bitmaps and downsampled to $8\times8$ by counting the
on-pixels in each non-overlapping $4\times4$ block. Each pixel is therefore an
integer in $[0, 16]$. A 1{,}797-image partition of that dataset is redistributed
with scikit-learn under a BSD-3-Clause licence.

This project commits a stratified random subset of 500 of those images --- 50
per digit --- drawn with \texttt{numpy.random.default\_rng(0)} and stored as
\texttt{data/digits\_sample.csv}: 64 pixel columns followed by the label. Pixel
values are copied through unchanged; there is no rescaling, augmentation, or
filtering. The subset exists so the file is about 75\,KB and can be shipped
with the project rather than downloaded. Because it is a subset, the accuracy
reported in Section \ref{sec:results} is lower than the full dataset gives, and
that is expected. The exact selection is reproduced in \texttt{data/README.md}.

Figure \ref{fig:01_digit_grid} shows one image per class, rendered directly
from the committed CSV.

\subsection{Model and evaluation}

The 500 images are split 70/30 into 350 training and 150 held-out images,
stratified by digit, with \texttt{random\_state=0}. A linear support vector
machine (\texttt{sklearn.svm.LinearSVC}, default regularisation, seeded) is fit
on the 350 training images using the 64 raw pixel values as features. No
feature extraction, dimensionality reduction, normalisation, or hyperparameter
search is performed --- the absence of those steps is the point of the example.

The model is evaluated once, on the 150 held-out images, by overall accuracy
and by the full $10\times10$ confusion matrix. There is no validation split and
no model selection, so no reuse of the held-out set occurs.

\subsection{Reproducibility}

Every number and both figures in this manuscript are produced by
\texttt{scripts/reproduce\_figures.py}, which reads only the committed CSV and
requires no network access. The subset draw, the train/test split, and the
solver are all seeded, so re-running the script reproduces the images shipped
with this project and the accuracy quoted below exactly:

\begin{verbatim}
python scripts/reproduce_figures.py
\end{verbatim}

%%%% EOF
